% =====================================================================
\chapter*{Conclusion}
\addcontentsline{toc}{chapter}{Conclusion}
\markboth{Conclusion}{Conclusion}
\label{chap-conclusion}
% =====================================================================

\section*{Rappel de la thèse}
\addcontentsline{toc}{section}{Rappel de la thèse}

Ce travail est parti d'une observation dont la simplicité est le principal
argument~: pour une page web, la segmentation dont un système de génération
augmentée par récupération visuelle a besoin existe déjà. Le moteur de mise en
page du navigateur a calculé la boîte englobante exacte de chaque élément avant
qu'un seul pixel ne soit recadré. PixelRAG~\citep{pixelrag2026} établit qu'un
modèle vision-langage peut récupérer directement sur des tuiles de pixels,
mais les découpe à l'aveugle~; MAGE-RAG~\citep{magerag2026} établit qu'un
contrôleur de graphe budgété surpasse une sélection statique, mais suppose les
éléments du document déjà constitués en amont de son propre pipeline. Les deux
laissent, chacun de son côté, la même pièce sur la table~: ce travail les
assemble par l'artefact que le premier jette au rendu et que le second n'a
jamais eu à construire --- le DOM.

Transformer cette observation en système a demandé quatre composants~: un tuilage
adaptatif aligné sur les frontières d'éléments, un graphe de preuves
multigranulaire s'étendant sur plusieurs documents, un lecteur vision-langage
adapté par contraste à l'évaluation de pertinence, et un contrôleur budgété
décidant à la requête quelles tuiles ouvrir. Le système est implémenté de bout en
bout, testé, et évalué sur des pages web réelles plutôt que sur un corpus figé.

\section*{Ce que le travail établit}
\addcontentsline{toc}{section}{Ce que le travail établit}

Deux résultats sont solidement étayés et constituent l'apport principal de ce
projet.

\paragraph{Le tuilage guidé par le DOM est meilleur et moins cher.} La fraction
des questions dont la réponse survit dans une unité récupérable passe de
\SI{26.4}{\percent} avec une grille aveugle à \SI{51.8}{\percent} avec nos
tuiles. Ce n'est pas un gain de classement mais un gain de \emph{possibilité}~:
une réponse absente de toute unité est perdue quelle que soit la qualité du
récupérateur en aval, de sorte que doubler ce taux double le nombre de questions
auxquelles le système peut du tout répondre.

Et ce gain s'obtient à un coût inférieur. Notre tuilage produit cinq fois plus
d'unités que la grille, mais chacune est treize fois plus petite~: le total des
pixels par page est \num{2.7} fois moindre. L'origine de ce résultat est
identifiée et formalisée précisément~--- la boîte \emph{serrée} d'une balise de
texte coulant, mesurée sur l'union des rectangles de ses lignes réellement
rendues plutôt que sur sa boîte de bloc CSS pleine largeur, élimine une aire de
fond blanc qu'un recadrage naïf aurait soumise en pure perte à l'encodeur
visuel ---~et elle en constitue l'enseignement technique le plus transférable.
Comme le nombre de tokens visuels croît avec l'aire de l'image dans la famille
de modèles utilisée, ce budget de pixels est un proxy direct du coût
d'inférence. Il n'y a donc pas d'arbitrage à faire~: la trouvabilité et le coût
s'améliorent ensemble.

\paragraph{L'adaptation contrastive du lecteur fonctionne nettement.} Sur un
partitionnement au niveau de l'article strictement disjoint, et en classant
chaque question contre toutes les tuiles de sa page plutôt que contre le vivier
miné, le Recall@1 passe de \num{0.433} à \num{0.785} et le MRR de \num{0.578} à
\num{0.856}.

La comparaison avec la référence lexicale porte un enseignement qui dépasse notre
système. Le modèle vision-langage non adapté obtient un MRR de \num{0.578},
pratiquement identique à celui d'une similarité cosinus TF-IDF (\num{0.579}), et
un Recall@1 légèrement inférieur --- mais un Recall@3 et un Recall@5 légèrement
supérieurs~: ce n'est pas un déficit général de récupération. Autrement dit~:
substituer un grand modèle multimodal à un scoreur lexical, sans l'entraîner à
la tâche de récupération, n'apporte pas d'avantage net sur le rang de tête, qui
est ce qui compte pour une sélection à un seul candidat. Tout le gain de
Recall@1 vient de l'adaptation contrastive spécifique.

\paragraph{Ces deux résultats partagent une même racine conceptuelle.} Ce
travail ne propose ni nouvelle architecture d'encodeur ni mécanisme
d'attention inédit~: le lecteur affiné reste un Qwen2-VL générique. Ce qui
change est \emph{où} passe la frontière de recadrage. En alignant la grille de
\anglais{patches} que consomme le modèle vision-langage sur les frontières
d'éléments du DOM plutôt que sur une géométrie arbitraire, le tuilage garantit
qu'une tuile porte un contenu visuel homogène~: c'est ce qui rend chaque pixel
soumis à l'encodeur pertinent pour la question posée, aussi bien au moment de
la récupération non entraînée qu'à celui de l'apprentissage contrastif, dont
l'objectif n'a plus à démêler plusieurs sujets mélangés dans une même séquence
de \anglais{patches}. Le gain de trouvabilité, la réduction du budget de
pixels et l'efficacité de l'adaptation contrastive sont ainsi trois
manifestations mesurables d'une seule décision de conception, prise avant que
l'encodeur ne voie la première tuile.

\section*{Le verrou identifié}
\addcontentsline{toc}{section}{Le verrou identifié}

Les composants en aval n'atteignent pas encore le niveau des deux précédents, et
le principal apport analytique de ce rapport est d'avoir localisé précisément
pourquoi.

Le contrôleur budgété n'établit pas de supériorité sur une sélection top-$k$
statique. Ce résultat a néanmoins produit un enchaînement méthodologique
satisfaisant~: nous avons formulé une hypothèse causale falsifiable --- la
frontière d'exploration ne franchissait jamais les limites de document, la
relation d'ordre de lecture ne reliant jamais deux pages ---, implémenté le
correctif déduit de ce diagnostic, et mesuré l'effet prédit. La couverture
inter-pages est passée de \SI{29.7}{\percent} à \SI{39.6}{\percent}, réduisant
l'écart avec la référence de dix points à un point. La parité est atteinte, la
supériorité non.

L'évaluation bout-en-bout a ensuite montré que le mode de défaillance dominant du
système assemblé est l'\emph{abstention explicite} --- quatorze prédictions sur
vingt --- et non la fabulation. Le lecteur n'échoue pas à interpréter les preuves
qu'on lui fournit~: il constate correctement qu'elles sont insuffisantes. La
défaillance est donc située dans la sélection, en amont.

Quatre observations convergent alors vers le même composant. Le contrôleur
n'atteint qu'une couverture complète sur un peu plus d'une question sur trois.
Douze des quatre-vingt-onze questions ont une couverture nulle, signe que le
score ne classe pas correctement même sur la page de départ. Le système assemblé
s'abstient parce que les preuves manquent. Et le tuilage structuré n'est pas
récompensé par une métrique aveugle à la structure --- ce qui explique le seul
résultat contraire à nos attentes, le découpage textuel plat devançant les tuiles
DOM sur le sous-ensemble strictement apparié.

Ce composant est le \textbf{score de pertinence du contrôleur}, qui reste un
substitut lexical alors même que l'évaluation du lecteur a validé un scoreur bien
supérieur sur exactement cette tâche. Le point d'extension qui permet la
substitution est déjà en place dans l'implémentation~; la seule barrière est le
coût de calcul, un passage avant du modèle par tuile candidate et par question
contre un score lexical essentiellement gratuit.

À ce premier verrou s'en ajoute un second, que le score ne recouvre pas. La
politique de parcours n'emprunte que trois des sept relations construites, et
l'adjacence de mise en page comme la hiérarchie de sections restent
inexploitées quel que soit le scoreur employé, la fonction de voisinage ne les
consultant pas. Le graphe est donc établi comme constructible depuis le DOM,
pas encore comme exploité.

Il faut dire cela clairement, car c'est une position favorable pour la suite. Le
système ne présente pas une défaillance diffuse dont on ignorerait l'origine,
mais deux composants validés et deux verrous nommés, localisés, instrumentés, et
dont les correctifs sont spécifiés séparément.

\section*{Enseignements méthodologiques}
\addcontentsline{toc}{section}{Enseignements méthodologiques}

Trois enseignements de méthode dépassent le cas particulier de ce système et
méritent d'être retenus.

\paragraph{Une métrique de récupération doit être invariante à la taille du
vivier.} Sur nos données, la grille aveugle atteint un Recall@5 de \num{0.991}
--- un score presque parfait qui n'exprime rien d'autre que le fait qu'elle ne
propose que cinq candidats par page. Rapporter des Recall bruts pour comparer des
stratégies de découpage produit une conclusion inversée. Le rang percentile, dont
l'espérance sous classement aléatoire vaut \num{0.5} quelle que soit la taille du
vivier, est la correction nécessaire. Cette précaution nous paraît sous-estimée
dans la littérature sur le découpage documentaire.

\paragraph{Le taux d'attachement mérite d'être mesuré explicitement.} La
proportion de questions dont la réponse survit à l'étape de découpage détermine
un plafond que le meilleur récupérateur du monde ne peut pas franchir, et c'est
sur cette métrique que l'écart entre les deux stratégies de tuilage est le plus
grand. Elle est pourtant absente des protocoles usuels, qui se concentrent sur la
qualité du classement en supposant l'attachement acquis.

\paragraph{Une évaluation par composant peut masquer un système qui ne
fonctionne pas.} Deux de nos quatre évaluations sont franchement positives, et le
système assemblé répond correctement à une question sur vingt. Sans l'évaluation
bout-en-bout, le rapport aurait pu conclure sur deux succès. C'est elle qui a
rendu lisible le mode de défaillance réel, et donc localisé le verrou. Nous
retenons de cette expérience qu'une évaluation par composant est nécessaire mais
jamais suffisante.

\section*{Perspectives}
\addcontentsline{toc}{section}{Perspectives}

La suite immédiate est déterminée par le diagnostic. Il s'agit de brancher le
lecteur validé dans le score de pertinence du contrôleur, à la place du
substitut lexical, et de rejouer à l'identique les évaluations du contrôleur et du
système assemblé. Une atténuation praticable du coût consiste à adopter
l'architecture classique du classement en deux étapes~: TF-IDF pour
présélectionner une courte liste, le lecteur affiné pour la réordonner. La même
substitution appliquée à l'évaluation du tuilage testerait directement si
l'avantage observé du texte plat est un artefact du scoreur lexical.

À plus longue portée, deux directions nous paraissent les plus prometteuses.
D'une part, faire réellement exploiter au contrôleur les relations du graphe qu'il
n'interroge pas encore~: remonter au titre de section d'une tuile ouverte, ou
atteindre l'en-tête d'un tableau dont une ligne a été retenue, sont des besoins de
récupération concrets auxquels le graphe répond déjà structurellement. D'autre
part, sortir du périmètre MediaWiki, ce qui permettrait de séparer empiriquement
ce qui relève de la liste de sélecteurs de nettoyage --- spécifique par
construction --- de ce qui relève de la logique de tuilage, qui ne l'est pas.

\section*{Mot de la fin}
\addcontentsline{toc}{section}{Mot de la fin}

L'affirmation centrale de ce travail résiste à l'examen~: le DOM contient bien la
segmentation que le RAG visuel reconstruit autrement, et l'exploiter rend
davantage de contenu trouvable pour un coût moindre. L'affirmation plus ambitieuse
--- qu'un graphe de preuves construit sur cette structure, parcouru par un
contrôleur budgété, produit un meilleur système de questions-réponses --- n'est pas
encore démontrée. Elle n'est pas réfutée non plus, car elle attend des
expériences identifiées, dont les obstacles sont un budget de calcul et une
politique de parcours à étendre, non une inconnue conceptuelle. Nous publions le pipeline complet, y compris les résultats qui ne
sont pas allés dans le sens attendu, parce que la localisation du verrou est
elle-même un résultat utilisable par qui reprendra ce travail.
